\chapter{Fundamentação teórica}
\lipsum[1-3]
\section{O paradigma LaTex}
Qual é a diferença entre WYSIWYG e WYSIWYM?


WYSIWYG: O paradigma WYSIWYG na pratica se refere a "Oque você vê é o que você pega", esse pega se refere a quando você estiver usando um editor de texto como Word ou Docs, você sempre fica clicando nos botoes para colocar em negrito, itálico, mudar a fonte, centralizar, aumentar o tamanho da fonte.  

    
WYSIWYM: Agora vamos explicar como funciona o paradigma do LaTex e porque ele funciona muito melhor, mesmo sendo mais complexo, o paradigma WISIWYM se refere a "O que você vê é oque você diz" pois no LaTex quando você precisa definir com códigos oque você quer dizer no seu texto ou modificar, um exemplo é: quando você quer colocar uma frase em negrito você usa o comando "textbf", ou outros códigos como "section" que serve para iniciar uma nova seção dentro de um paragrafo.

    
No começo parece complexo mas quando você começa a entender como funciona o código fica bem mais fácil usar o LaTex que um editor de texto, pois o LaTex coloca seu texto em ABNT automaticamente, deixa ele todo organizado, faz resumo, isso tudo com pequenas linhas de código, oque no Word demoraria horas, talvez dias, tempo que você poderia gastar escrevendo mais ou fazendo varias outras coisas.

\subsection{Os diferentes modelos para criação}
No LaTex na primeira linha do código, você precisa escolher qual tipo de modelo você vai utilizar no seu texto existe varias opções e cada um serve para uma coisa em especifico, vou citar alguns deles, existe o report, beamer e article, essas são as classes de documento padrões do LaTex.

Beamer: A classe de documento beamer cria uma forma de apresentação para você, ela é bastante usada para apresentações acadêmicas, seminários ou palestras, ela permite inserir tabelas, formulas, imagens entre outras coisas na sua apresentação.


Article: A classe article e uma das mais utilizadas no LaTex, ela serve para criar documentos menores, ela é muito utilizada para criar artigos científicos, resumos, listas de exercício e relatórios simples, devido a sua forma de estrutura ela permite organizar o texto em seções e subseções, deixando o texto com muito mais coerência para ler.


Report: A classe report funciona de uma forma parecida com a classe article, porem a classe report serve para textos maiores como teses ou dissertações, na classe report também podemos utilizar seções e subseções, porem podemos também utilizar subsubseções e capítulos, diferenciando da article. 


\subsection{Os diferentes tamanhos de texto}
No LaTex quando você quer deixar alguma parte do texto tem varios tamanhos ja predefinidos para você utilizar no seu texto, tem tamanhos desde o mais pequeno ate tamanhos grotescamente grandes, para utilizar esse tamanhos de teste basta apenas digitar qual o nome do tamanho que voce que e entre colchetes colocar o texto que voce quer que cresça, vou deixar amostra alguns do principais tamanhos:

\begin{itemize}
    \item \textbackslash Tiny (Tamanho minusculo)
    \item \textbackslash small (Tamanho pequeno)
    \item \textbackslash large (Tamanho grande)
    \item \textbackslash huge (Tamanho gigante)
\end{itemize}


Tem muito mais tamanhos para se utilizar, porem estes sao os mais utilizados, durante o texto quando voce quer deixar algo mais visivel, ou deixar algo menor os tamanhos sao muito uteis, um exemplo de utilidade, é quando voce vai colocar um codigo no seu artigo, geralmente as pessoas colocam em uma fonte menor que a normal por ter muito codigo na maioria das vezes. 













